\documentclass[10pt,letterpaper]{article}
\usepackage{cloudsuite}
\renewcommand{\DocShort}{CH. 10}
\begin{document}
\coverpage{Chapter 10: Cloud Security and Cybersecurity Access-Oriented Mechanisms}{Cornell study notes and applied review | MECHANISMS}{Concepts \textbullet{} decisions \textbullet{} verification artifacts}
\sourceStatement
\tableofcontents
\newpage
\section{Learning contract}
\begin{takeaway}
Design a defense-in-depth access control chain.
\end{takeaway}
\textbf{Chapter scope:} Cryptography, identity, network controls, testing and monitoring.\\
\sourceRef{10}{269}
\subsection{Mapped source sections}
\begin{itemize}
\item 10.1 Encryption
\item 10.2 Hashing
\item 10.3 Digital Signature
\item 10.4 Cloud-Based Security Groups
\item 10.5 Public Key Infrastructure (PKI) System
\item 10.6 Single Sign-On (SSO) System
\item 10.7 Hardened Virtual Server Image
\item 10.8 Firewall
\item 10.9 Virtual Private Network (VPN)
\item 10.10 Biometric Scanner
\item 10.11 Multi-Factor Authentication (MFA) System
\item 10.12 Identity and Access Management (IAM) System
\item 10.13 Intrusion Detection System (IDS)
\item 10.14 Penetration Testing Tool
\item 10.15 User Behavior Analytics (UBA) System
\item 10.16 Third-Party Software Update Utility
\item 10.17 Network Intrusion Monitor
\item 10.18 Authentication Log Monitor
\item 10.19 VPN Monitor
\item 10.20 Additional Cloud Security Access-Oriented Practices and Technologies
\end{itemize}
\section{Cornell cue and notes matrix}
\small\setlength{\LTleft}{0pt}\setlength{\LTright}{0pt}
\begin{longtable}{>{\RaggedRight\arraybackslash}p{0.29\textwidth}>{\RaggedRight\arraybackslash}p{0.63\textwidth}}
\toprule\textbf{Cue / essential question}&\textbf{Explanatory notes and reasoning}\\\midrule\endfirsthead\toprule\textbf{Cue}&\textbf{Notes}\\\midrule\endhead
\textbf{What separates encryption, hashing and signatures?} & Encryption protects confidentiality with keys; hashing provides a digest; digital signatures connect integrity verification with a signing key. None replaces access governance.\\[3pt]\textcolor{petrol}{\emph{Practice:} Assign the correct primitive to three scenarios.}\\[8pt]\midrule
\textbf{What does a PKI system contribute?} & Certificate issuance and trust chains bind public keys to identities or services. Expiry and revocation are lifecycle controls, not cosmetic details.\\[3pt]\textcolor{petrol}{\emph{Practice:} Document certificate rotation and failure handling.}\\[8pt]\midrule
\textbf{How do IAM, SSO and MFA interact?} & IAM establishes identity and authorization policy; SSO shares an authentication experience; MFA strengthens authentication assurance.\\[3pt]\textcolor{petrol}{\emph{Practice:} Trace authorization after user login and token issuance.}\\[8pt]\midrule
\textbf{What layers enforce network access?} & Security groups, firewalls and VPNs filter or protect different portions of communication. Define the boundary and rule owner for each layer.\\[3pt]\textcolor{petrol}{\emph{Practice:} Trace one permitted and one prohibited path.}\\[8pt]\midrule
\textbf{Why harden server images and update software?} & Repeatable hardened baselines and ongoing updates reduce exploitable conditions across a fleet.\\[3pt]\textcolor{petrol}{\emph{Practice:} Define baseline controls and drift tests.}\\[8pt]\midrule
\textbf{How do detection and testing differ?} & IDS, network/authentication/VPN monitoring and UBA observe activity; penetration testing challenges assumptions under authorization.\\[3pt]\textcolor{petrol}{\emph{Practice:} Relate one detection signal to a controlled security test.}\\[8pt]\midrule
\end{longtable}\normalsize
\section{Synthesis and architecture connections}
Treat each mechanism as a deployable capability, not a complete architecture. Identify the infrastructure prerequisite, management-plane owner, dependent controls and the negative/failure evidence expected from an implementation.
\begin{workbox}[title={Applied deliverable}]
Design a defense-in-depth access control chain.\\[4pt]Create a one-page architecture or operations artifact with: (1) assumptions, (2) a decision and alternatives, (3) control ownership, (4) observable acceptance criteria, and (5) a failed-case test.
\end{workbox}
\section{Active recall}
\begin{enumerate}
\item What separates encryption, hashing and signatures?
\item What does a PKI system contribute?
\item How do IAM, SSO and MFA interact?
\item What layers enforce network access?
\end{enumerate}
\subsection{Answer-check rubric}
\begin{itemize}
\item A correct response distinguishes the named concepts rather than relying on vendor names.
\item The explanation identifies an operational consequence, boundary or tradeoff.
\item At least one measurement, test or documentary artifact could demonstrate the claim.
\end{itemize}
\section{Blank Cornell capture sheet}
\textbf{Topic / source subsection:}\enspace\rule{0.65\linewidth}{0.4pt}\\[6pt]
\begin{longtable}{p{0.28\textwidth}p{0.64\textwidth}}
\toprule Cues / questions & Notes / evidence \\ \midrule
\rule{0pt}{42pt} & \\ \midrule
\rule{0pt}{42pt} & \\ \midrule
\rule{0pt}{42pt} & \\ \midrule
\rule{0pt}{42pt} & \\ \midrule
\rule{0pt}{42pt} & \\ \midrule
\end{longtable}
\textbf{Summary in three sentences:}\\[16pt]\hrulefill\\[16pt]\hrulefill\\[16pt]\hrulefill

\refs
\end{document}
